% =====================================================================
\chapter{Matrices y aplicaciones lineales}
\label{tema:3}
% =====================================================================

Una matriz es una tabla de números, pero su papel en el álgebra lineal es el de una función: transforma vectores en vectores respetando las combinaciones lineales. Este tema establece esa correspondencia entre matrices y aplicaciones lineales. Se sigue la sección~2.1 de Strang~\cite{strang} y la sesión «The Geometry of Linear Equations» de 18.06SC~\cite{mit1806sc}; las definiciones proceden de~\cite[§2.2 y §2.7]{mml} y de~\cite[cap.~6 y §8.1]{vmls}.

\section{Matrices}

\begin{definicion}[Matriz]
Sean $m,n\ge 1$. Una \emph{matriz} real de tamaño $m\times n$ es una tabla $A=(A_{ij})$ de $mn$ números reales dispuestos en $m$ filas y $n$ columnas; $A_{ij}$ es la entrada de la fila $i$ y la columna $j$. El conjunto de estas matrices se denota $\R^{m\times n}$. Se escribe $a_j\in\R^{m}$ para la columna $j$ y $A=[\,a_1\ \cdots\ a_n\,]$.
\end{definicion}

Un vector de $\R^{d}$ se identifica con una matriz $d\times 1$ (vector columna). Las matrices del mismo tamaño se suman y se multiplican por escalares entrada a entrada, y con esas operaciones $\R^{m\times n}$ es un espacio vectorial.

\begin{definicion}[Traspuesta e identidad]
La \emph{traspuesta} de $A\in\R^{m\times n}$ es la matriz $A^{\top}\in\R^{n\times m}$ con $(A^{\top})_{ij}=A_{ji}$: las filas de $A$ son las columnas de $A^{\top}$. Una matriz cuadrada es \emph{simétrica} si $A^{\top}=A$. La \emph{identidad} $I_n\in\R^{n\times n}$ tiene unos en la diagonal y ceros fuera de ella; sus columnas son los vectores $e_1,\dots,e_n$ de la base canónica.
\end{definicion}

Si $u\in\R^{d}$ es una columna, $u^{\top}$ es una matriz $1\times d$ (vector fila). Esta notación permite escribir el producto escalar como $u\cdot v=u^{\top}v$.

\section{El producto de una matriz por un vector}

\begin{definicion}[Producto matriz-vector]\label{def:matvec}
Sean $A\in\R^{m\times n}$ y $x\in\R^{n}$. El producto $Ax\in\R^{m}$ es el vector de componentes
\[
(Ax)_i=\sum_{j=1}^{n}A_{ij}x_j,\qquad i=1,\dots,m.
\]
El producto solo está definido si el número de columnas de $A$ coincide con el número de componentes de $x$.
\end{definicion}

\begin{proposicion}[Dos lecturas del producto]\label{prop:filas-columnas}
Sea $A\in\R^{m\times n}$ con filas $r_1^{\top},\dots,r_m^{\top}$ y columnas $a_1,\dots,a_n$. Para todo $x\in\R^{n}$:
\begin{enumerate}
  \item (por filas) la componente $i$ de $Ax$ es el producto escalar $r_i\cdot x$;
  \item (por columnas) $Ax=x_1a_1+x_2a_2+\dots+x_na_n$ es una combinación lineal de las columnas de $A$.
\end{enumerate}
\end{proposicion}

\addtocounter{definicion}{1}%

Las dos lecturas corresponden a las dos imágenes de un sistema de ecuaciones lineales $Ax=b$ que Strang llama imagen por filas e imagen por columnas~\cite[§2.1]{strang}: por filas, cada ecuación es un hiperplano y la solución es su intersección; por columnas, se buscan los coeficientes con los que las columnas de $A$ se combinan para dar $b$.

\section{Aplicaciones lineales}

\begin{definicion}[Aplicación lineal]\label{def:lineal}
Una aplicación $T\colon\R^{n}\to\R^{m}$ es \emph{lineal} si
\[
T(au+bv)=a\,T(u)+b\,T(v)\qquad\text{para todos } u,v\in\R^{n},\ a,b\in\R.
\]
\end{definicion}

Tomando $a=b=0$ se obtiene $T(0)=0$: una aplicación lineal fija el origen. Por inducción, $T$ transforma cualquier combinación lineal en la combinación lineal de las imágenes con los mismos coeficientes.

\begin{proposicion}\label{prop:matriz-lineal}
Para toda $A\in\R^{m\times n}$, la aplicación $x\mapsto Ax$ es lineal.
\end{proposicion}

El resultado recíproco es el que justifica la frase «una matriz es una aplicación lineal».

\begin{teorema}[Matriz de una aplicación lineal]\label{teo:representacion}
Sea $T\colon\R^{n}\to\R^{m}$ lineal y sea $A\in\R^{m\times n}$ la matriz cuyas columnas son $T(e_1),\dots,T(e_n)$. Entonces $T(x)=Ax$ para todo $x\in\R^{n}$, y $A$ es la única matriz con esta propiedad.
\end{teorema}

El \cref{teo:representacion} da una receta práctica: para conocer una aplicación lineal basta saber adónde envía los vectores de la base canónica, y esas imágenes son las columnas de su matriz. El \cref{tab:transformaciones} recoge las transformaciones del plano más habituales, y la \cref{fig:cizalla} muestra el efecto de una cizalla sobre el cuadrado unidad.

\begin{table}[h]
\centering\small
\caption{Transformaciones lineales del plano. Las columnas de cada matriz son las imágenes de $e_1$ y $e_2$.}\label{tab:transformaciones}
\begin{tabular}{@{}lccc@{}}
\toprule
Transformación & Matriz & $e_1\mapsto$ & $e_2\mapsto$ \\
\midrule
Rotación de ángulo $\theta$ & $R_\theta=\begin{pmatrix}\cos\theta&-\sin\theta\\ \sin\theta&\cos\theta\end{pmatrix}$ & $(\cos\theta,\sin\theta)$ & $(-\sin\theta,\cos\theta)$ \\[10pt]
Rotación de $90^\circ$ & $\begin{pmatrix}0&-1\\1&0\end{pmatrix}$ & $(0,1)$ & $(-1,0)$ \\[8pt]
Escalado & $\begin{pmatrix}2&0\\0&0{,}5\end{pmatrix}$ & $(2,0)$ & $(0;\,0{,}5)$ \\[8pt]
Cizalla & $\begin{pmatrix}1&1\\0&1\end{pmatrix}$ & $(1,0)$ & $(1,1)$ \\[8pt]
Simetría respecto de $y=x$ & $\begin{pmatrix}0&1\\1&0\end{pmatrix}$ & $(0,1)$ & $(1,0)$ \\[8pt]
Proyección sobre el eje $x$ & $\begin{pmatrix}1&0\\0&0\end{pmatrix}$ & $(1,0)$ & $(0,0)$ \\
\bottomrule
\end{tabular}
\end{table}

\addtocounter{figure}{1}%

\begin{proposicion}[Las rotaciones conservan la norma]\label{prop:rotacion}
Para todo $\theta\in\R$ y todo $x\in\R^{2}$, $\norm{R_\theta x}=\norm{x}$.
\end{proposicion}

Las rotaciones conservan también los productos escalares y, por tanto, los ángulos (ejercicio~\ref{eje:4-ortogonal}). Los escalados no: con la matriz de escalado del \cref{tab:transformaciones}, $e_1+e_2$, de norma $\sqrt2$, pasa a $(2;\,0{,}5)$, de norma $\sqrt{4{,}25}$.

\begin{observacion}[Matrices en los modelos de lenguaje]\label{obs:qkv}
Casi todo el cálculo de un Transformer consiste en aplicar matrices aprendidas a vectores. En la atención (módulo~13), cada embedding $x\in\R^{d}$ se transforma en una \emph{consulta} $q=W_Qx$, una \emph{clave} $k=W_Kx$ y un \emph{valor} $v=W_Vx$, con $W_Q,W_K,W_V$ matrices $d_k\times d$. Estas matrices seleccionan qué rasgos del embedding intervienen en cada comparación.
\end{observacion}

\begin{observacion}[Matrices de transición]
La matriz de transición $P$ de un modelo de bigramas (módulo~1, tema~2) es estocástica por filas. Si $\pi_t\in\R^{K}$ es la distribución del token en la posición $t$, la de la posición siguiente es $\pi_{t+1}=P^{\top}\pi_t$, cuya componente $j$ es $\sum_iP_{ij}(\pi_t)_i$. La matriz $P^{\top}$ tiene columnas que suman~1; es la convención que usa Strang para las \emph{matrices de Markov} (ejercicio~\ref{eje:3-markov}).
\end{observacion}

